\subsection{Biểu đồ use case phân rã Quản lý flashcard}
\label{subsection:2.2.2}

\begin{figure}[H]
    \centering
    \includegraphics[width=0.9\textwidth]{Hinhve/UC/UseCase_QuanLyFlashcard.png}
    \caption{Biểu đồ use case phân rã: Quản lý flashcard}
    \label{fig:usecase_quanlyflashcard}
\end{figure}

\textbf{Biểu đồ \ref{fig:usecase_quanlyflashcard}} mô tả cách học sinh thực hiện các thao tác quản lý flashcard trên hệ thống. Thông qua Use Case chính là "Quản lý flashcard", học sinh có thể tạo, chỉnh sửa, xóa bộ flashcard và các flashcard cá nhân, đồng thời quản lý chúng theo category và thực hiện các thao tác sao chép, import/export. Các chức năng này cho phép học sinh tổ chức hiệu quả nội dung học tập của mình, đảm bảo tính linh hoạt và dễ dàng trong việc quản lý tài liệu học tập.

\subsection{Biểu đồ use case phân rã Học flashcard}
\label{subsection:2.2.3}

\begin{figure}[H]
    \centering
    \includegraphics[width=0.9\textwidth]{Hinhve/UC/UseCase_HocFlashcard.png}
    \caption{Biểu đồ use case phân rã: Học flashcard}
    \label{fig:usecase_hocflashcard}
\end{figure}

\textbf{Biểu đồ \ref{fig:usecase_hocflashcard}} mô tả cách học sinh thực hiện quá trình học tập với flashcard trên hệ thống. Thông qua Use Case chính là "Học flashcard", học sinh có thể bắt đầu phiên học, học theo nhiều chế độ khác nhau (lặp lại, kiểm tra nhanh, ngẫu nhiên), đánh giá mức độ nhớ, và tương tác với AI để được hỗ trợ học tập. Các chức năng này cho phép học sinh học tập hiệu quả với nhiều phương thức khác nhau, tận dụng công nghệ AI để nâng cao trải nghiệm học tập và củng cố kiến thức.

\subsection{Biểu đồ use case phân rã Làm kiểm tra}
\label{subsection:2.2.4}

\begin{figure}[H]
    \centering
    \includegraphics[width=0.9\textwidth]{Hinhve/UC/UseCase_LamKiemTra.png}
    \caption{Biểu đồ use case phân rã: Làm kiểm tra}
    \label{fig:usecase_lamkiemtra}
\end{figure}

\textbf{Biểu đồ \ref{fig:usecase_lamkiemtra}} mô tả cách học sinh thực hiện các thao tác làm bài kiểm tra trên hệ thống. Thông qua Use Case chính là "Làm kiểm tra", học sinh có thể xem danh sách bài kiểm tra, bắt đầu làm bài, trả lời các loại câu hỏi khác nhau (trắc nghiệm, điền từ, tự luận), nộp bài và xem kết quả. Các chức năng này cho phép học sinh đánh giá kiến thức một cách toàn diện, theo dõi tiến độ học tập và cải thiện hiệu quả học tập thông qua việc làm lại bài kiểm tra khi được phép.

\subsection{Biểu đồ use case phân rã Quản lý lịch sử kiểm tra}
\label{subsection:2.2.5}

\begin{figure}[H]
    \centering
    \includegraphics[width=0.9\textwidth]{Hinhve/UC/UseCase_QuanLyLichSuKiemTra.png}
    \caption{Biểu đồ use case phân rã: Quản lý lịch sử kiểm tra}
    \label{fig:usecase_quanlylichsukiemtra}
\end{figure}

\textbf{Biểu đồ \ref{fig:usecase_quanlylichsukiemtra}} mô tả cách học sinh quản lý và xem lại lịch sử kiểm tra trên hệ thống. Thông qua Use Case chính là "Quản lý lịch sử kiểm tra", học sinh có thể xem danh sách và chi tiết kết quả các bài kiểm tra đã làm, lọc lịch sử theo nhiều tiêu chí khác nhau, và xuất báo cáo kết quả. Các chức năng này cho phép học sinh theo dõi tiến độ học tập một cách chi tiết, phân tích hiệu suất học tập và đưa ra các quyết định cải thiện phù hợp.

\subsection{Biểu đồ use case phân rã Quản lý lớp học}
\label{subsection:2.2.6}

\begin{figure}[H]
    \centering
    \includegraphics[width=0.9\textwidth]{Hinhve/UC/UseCase_QuanLyLopHoc.png}
    \caption{Biểu đồ use case phân rã: Quản lý lớp học}
    \label{fig:usecase_quanlylophoc}
\end{figure}

\textbf{Biểu đồ \ref{fig:usecase_quanlylophoc}} mô tả cách giáo viên thực hiện các thao tác quản lý lớp học trên hệ thống. Thông qua Use Case chính là "Quản lý lớp học", giáo viên có thể tạo, chỉnh sửa, xóa lớp học, quản lý thành viên (mời, chấp nhận, từ chối, xóa), quản lý nội dung học tập và tạo mã tham gia lớp. Các chức năng này cho phép giáo viên kiểm soát hiệu quả hoạt động của lớp học, quản lý học sinh và nội dung giảng dạy một cách có tổ chức, đảm bảo môi trường học tập tốt nhất cho học sinh.

\subsection{Biểu đồ use case phân rã Quản lý bài kiểm tra}
\label{subsection:2.2.7}

\begin{figure}[H]
    \centering
    \includegraphics[width=0.9\textwidth]{Hinhve/UC/UseCase_QuanLyBaiKiemTra.png}
    \caption{Biểu đồ use case phân rã: Quản lý bài kiểm tra}
    \label{fig:usecase_quanlybaikiemtra}
\end{figure}

\textbf{Biểu đồ \ref{fig:usecase_quanlybaikiemtra}} mô tả cách giáo viên thực hiện các thao tác quản lý bài kiểm tra trên hệ thống. Thông qua Use Case chính là "Quản lý bài kiểm tra", giáo viên có thể tạo, chỉnh sửa, xóa bài kiểm tra, cấu hình các thông số (loại câu hỏi, số lượng, thời gian, hiển thị đáp án), xem kết quả làm bài của học sinh và xử lý yêu cầu làm lại bài. Các chức năng này cho phép giáo viên tạo và quản lý bài kiểm tra một cách linh hoạt, đánh giá chính xác năng lực học sinh và hỗ trợ học sinh cải thiện kết quả học tập.

\subsection{Biểu đồ use case phân rã Quản lý nội dung học tập}
\label{subsection:2.2.8}

\begin{figure}[H]
    \centering
    \includegraphics[width=0.9\textwidth]{Hinhve/UC/UseCase_QuanLyNoiDungHocTap.png}
    \caption{Biểu đồ use case phân rã: Quản lý nội dung học tập}
    \label{fig:usecase_quanlynoidunghoctap}
\end{figure}

\textbf{Biểu đồ \ref{fig:usecase_quanlynoidunghoctap}} mô tả cách giáo viên thực hiện các thao tác quản lý nội dung học tập trên hệ thống. Thông qua Use Case chính là "Quản lý nội dung học tập", giáo viên có thể tạo, chỉnh sửa, xóa bộ flashcard và tài liệu học tập, gán hoặc hủy gán nội dung cho lớp học, và phân loại nội dung theo chủ đề. Các chức năng này cho phép giáo viên tổ chức và quản lý tài nguyên học tập một cách hiệu quả, đảm bảo học sinh có đầy đủ tài liệu cần thiết cho việc học tập.

\subsection{Biểu đồ use case phân rã Xem thống kê lớp học}
\label{subsection:2.2.9}

\begin{figure}[H]
    \centering
    \includegraphics[width=0.9\textwidth]{Hinhve/UC/UseCase_XemThongKeLopHoc.png}
    \caption{Biểu đồ use case phân rã: Xem thống kê lớp học}
    \label{fig:usecase_xemthongkelophoc}
\end{figure}

\textbf{Biểu đồ \ref{fig:usecase_xemthongkelophoc}} mô tả cách giáo viên xem và phân tích thống kê lớp học trên hệ thống. Thông qua Use Case chính là "Xem thống kê lớp học", giáo viên có thể xem thống kê tổng quan, tiến độ học tập của từng học sinh, tỉ lệ nhớ/quên, số flashcard đã hoàn thành, mức độ tiến bộ theo thời gian và kết quả bài kiểm tra. Các chức năng này cho phép giáo viên theo dõi và đánh giá hiệu quả học tập của lớp học một cách chi tiết, từ đó đưa ra các biện pháp hỗ trợ và cải thiện phù hợp cho học sinh.

\subsection{Biểu đồ use case phân rã Tham gia vào lớp học}
\label{subsection:2.2.10}

\begin{figure}[H]
    \centering
    \includegraphics[width=0.9\textwidth]{Hinhve/UC/UseCase_ThamGiaVaoLopHoc.png}
    \caption{Biểu đồ use case phân rã: Tham gia vào lớp học}
    \label{fig:usecase_thamgialophoc}
\end{figure}

\textbf{Biểu đồ \ref{fig:usecase_thamgialophoc}} mô tả cách học sinh tham gia và tương tác với lớp học trên hệ thống. Thông qua Use Case chính là "Tham gia vào lớp học", học sinh có thể tìm kiếm lớp học công khai, tham gia bằng mã lớp, gửi yêu cầu tham gia, xem danh sách và chi tiết lớp học, rời khỏi lớp học, và xem các tài nguyên học tập trong lớp. Các chức năng này cho phép học sinh dễ dàng tham gia vào các lớp học phù hợp, tiếp cận nội dung học tập và bài kiểm tra một cách thuận tiện.

\subsection{Biểu đồ use case phân rã Xem tiến độ học tập}
\label{subsection:2.2.11}

\begin{figure}[H]
    \centering
    \includegraphics[width=0.9\textwidth]{Hinhve/UC/UseCase_XemTienDoHocTap.png}
    \caption{Biểu đồ use case phân rã: Xem tiến độ học tập}
    \label{fig:usecase_xemtiendohoctap}
\end{figure}

\textbf{Biểu đồ \ref{fig:usecase_xemtiendohoctap}} mô tả cách học sinh xem và theo dõi tiến độ học tập trên hệ thống. Thông qua Use Case chính là "Xem tiến độ học tập", học sinh có thể xem tổng quan tiến độ, số flashcard đã học, tỉ lệ nhớ/quên, thời gian học tập, tiến độ theo từng bộ flashcard, biểu đồ tiến độ theo thời gian và danh sách flashcard cần ôn lại. Các chức năng này cho phép học sinh tự đánh giá hiệu quả học tập của mình, xác định các điểm cần cải thiện và lập kế hoạch học tập phù hợp.

\subsection{Biểu đồ use case phân rã Quản lý category}
\label{subsection:2.2.12}

\begin{figure}[H]
    \centering
    \includegraphics[width=0.9\textwidth]{Hinhve/UC/UseCase_QuanLyCategory.png}
    \caption{Biểu đồ use case phân rã: Quản lý category}
    \label{fig:usecase_quanlycategory}
\end{figure}

\textbf{Biểu đồ \ref{fig:usecase_quanlycategory}} mô tả cách học sinh quản lý danh mục (category) để phân loại flashcard trên hệ thống. Thông qua Use Case chính là "Quản lý category", học sinh có thể tạo, chỉnh sửa, xóa category, gán hoặc xóa flashcard khỏi category, xem danh sách flashcard theo category và sắp xếp category. Các chức năng này cho phép học sinh tổ chức flashcard một cách có hệ thống, dễ dàng tìm kiếm và quản lý nội dung học tập theo các chủ đề khác nhau.

\subsection{Biểu đồ use case phân rã Quản lý người dùng}
\label{subsection:2.2.13}

\begin{figure}[H]
    \centering
    \includegraphics[width=0.9\textwidth]{Hinhve/UC/UseCase_QuanLyNguoiDung.png}
    \caption{Biểu đồ use case phân rã: Quản lý người dùng}
    \label{fig:usecase_quanlynguoidung}
\end{figure}

\textbf{Biểu đồ \ref{fig:usecase_quanlynguoidung}} mô tả cách quản trị viên thực hiện các thao tác quản trị người dùng trên hệ thống. Thông qua Use Case chính là "Quản lý người dùng", Admin có thể xem danh sách, tạo mới, chỉnh sửa và xóa người dùng, thay đổi vai trò, khóa/mở khóa tài khoản, xem chi tiết hoạt động, tìm kiếm và lọc người dùng. Các chức năng này cho phép quản trị viên kiểm soát hiệu quả cơ sở dữ liệu người dùng, đảm bảo tính chính xác, bảo mật và cập nhật thông tin kịp thời.

\subsection{Biểu đồ use case phân rã Xem thống kê toàn hệ thống}
\label{subsection:2.2.14}

\begin{figure}[H]
    \centering
    \includegraphics[width=0.9\textwidth]{Hinhve/UC/UseCase_XemThongKeToanHeThong.png}
    \caption{Biểu đồ use case phân rã: Xem thống kê toàn hệ thống}
    \label{fig:usecase_xemthongketoanhethong}
\end{figure}

\textbf{Biểu đồ \ref{fig:usecase_xemthongketoanhethong}} mô tả cách quản trị viên xem và phân tích thống kê toàn hệ thống. Thông qua Use Case chính là "Xem thống kê toàn hệ thống", Admin có thể xem tổng quan hệ thống, thống kê số lượng người dùng, hoạt động học tập, hiệu suất sử dụng hệ thống, thống kê theo vai trò, theo khoảng thời gian, thống kê lớp học và flashcard. Các chức năng này cho phép quản trị viên theo dõi và đánh giá hiệu quả hoạt động của toàn bộ hệ thống, từ đó đưa ra các quyết định quản lý và cải thiện hệ thống phù hợp.

\subsection{Biểu đồ use case phân rã Quản lý thông tin cá nhân}
\label{subsection:2.2.15}

\begin{figure}[H]
    \centering
    \includegraphics[width=0.9\textwidth]{Hinhve/UC/UseCase_QuanLyThongTinCaNhan.png}
    \caption{Biểu đồ use case phân rã: Quản lý thông tin cá nhân}
    \label{fig:usecase_quanlythongtincanhan}
\end{figure}

\textbf{Biểu đồ \ref{fig:usecase_quanlythongtincanhan}} mô tả cách người dùng (học sinh và giáo viên) quản lý thông tin cá nhân trên hệ thống. Thông qua Use Case chính là "Quản lý thông tin người dùng", người dùng có thể xem, cập nhật thông tin cá nhân, thay đổi mật khẩu, thông tin liên hệ và xem lịch sử hoạt động cá nhân. Các chức năng này cho phép người dùng duy trì thông tin tài khoản chính xác và cập nhật, đảm bảo tính bảo mật và theo dõi hoạt động học tập của bản thân.
